\chapter{Introducción a la topología general}
\label{ch:topologia-general}
\index{Espacio topológico}\index{Subbase}\index{Axiomas de separación}\index{Producto de espacios}\index{Teorema de Tychonoff}

\section{Espacios topológicos}

La topología general nace de la abstracción de las propiedades de los espacios métricos que no dependen de la distancia explícita. Hausdorff, en su \emph{Grundzüge der Mengenlehre} (1914), formalizó esta idea definiendo un espacio topológico mediante una familia de conjuntos abiertos. Este salto conceptual permitió estudiar espacios que no provienen de ninguna métrica: topologías débiles en análisis funcional, topologías de Zariski en geometría algebraica, y topologías de orden en teoría de conjuntos.

\begin{definition}
\label{def:espacio-topologico}
Un \emph{espacio topológico} es un par $(X,\tau)$ donde $X$ es un conjunto y $\tau\subseteq\mathcal{P}(X)$ es una familia, llamada \emph{topología}, que satisface:
\begin{enumerate}
  \item[(T1)] $\emptyset,X\in\tau$.
  \item[(T2)] La unión arbitraria de elementos de $\tau$ pertenece a $\tau$.
  \item[(T3)] La intersección finita de elementos de $\tau$ pertenece a $\tau$.
\end{enumerate}
\end{definition}

\begin{example}
\label{ej:metrico-topologico}
Todo espacio métrico induce una topología métrica (\cref{teo:topologia-metrica}).
\end{example}

\begin{example}
\label{ej:cofinita}
La \emph{topología cofinita} sobre un conjunto infinito $X$, donde los abiertos son $\emptyset$ y los complementos de conjuntos finitos, no es metrizable.
\end{example}

\section{Bases y subbases}

\begin{definition}
\label{def:base-topologia}
Una \emph{base} de una topología $\tau$ es una subfamilia $\mathcal{B}\subseteq\tau$ tal que todo abierto no vacío es unión de elementos de $\mathcal{B}$.
\end{definition}

\begin{definition}
\label{def:subbase}
Una \emph{subbase} es una familia $\mathcal{S}$ tal que las intersecciones finitas de sus elementos forman una base.
\end{definition}

\section{Continuidad}

\begin{definition}
\label{def:continuidad-topologica}
Una función $f\colon(X,\tau_X)\to(Y,\tau_Y)$ es \emph{continua} si $f^{-1}(V)\in\tau_X$ para todo $V\in\tau_Y$.
\end{definition}

\begin{proposition}
\label{prop:comp-topologica}
La composición de funciones continuas es continua.
\end{proposition}

\section{Compacidad}

\begin{definition}
\label{def:compacto-topologia}
Un espacio topológico $X$ es \emph{compacto} si de todo recubrimiento abierto se puede extraer un subrecubrimiento finito.
\end{definition}

\begin{theorem}
\label{teo:compacto-cerrado}
En un espacio compacto Hausdorff, los compactos coinciden con los cerrados.
\end{theorem}

\section{Conexidad}

\begin{definition}
\label{def:conexo-topologia}
Un espacio topológico es \emph{conexo} si no puede escribirse como unión de dos abiertos disjuntos no vacíos.
\end{definition}

\begin{proposition}
\label{prop:imagen-conexa}
La imagen continua de un conexo es conexa.
\end{proposition}

\section{Axiomas de separación}

\begin{definition}
\label{def:separacion}
\begin{itemize}
  \item $T_0$: para puntos distintos, al menos uno tiene un entorno que no contiene al otro.
  \item $T_1$: para puntos distintos, cada uno tiene un entorno que no contiene al otro.
  \item $T_2$ (Hausdorff): puntos distintos tienen entornos disjuntos.
  \item $T_3$ (regular): cerrado y punto fuera de él tienen entornos disjuntos.
  \item $T_4$ (normal): dos cerrados disjuntos tienen entornos disjuntos.
\end{itemize}
\end{definition}

\begin{theorem}[Urysohn]
\label{teo:urysohn-general}
En un espacio normal, cualesquiera dos cerrados disjuntos pueden separarse por una función continua $f\colon X\to[0,1]$.
\end{theorem}

\begin{theorem}[Tietze]
\label{teo:tietze-general}
En un espacio normal, toda función continua definida en un cerrado y con valores en $\mathbb{R}$ puede extenderse continuamente a todo $X$.
\end{theorem}

\section{Productos de espacios}

\begin{definition}
\label{def:producto-topologico}
La \emph{topología producto} sobre $\prod_{i\in I}X_i$ es la generada por las proyecciones $\pi_j$.
\end{definition}

\section{Teorema de Tychonoff (panorama)}

\begin{theorem}[Tychonoff]
\label{teo:tychonoff}
El producto arbitrario de espacios compactos es compacto en la topología producto.
\end{theorem}

Este teorema, uno de los más profundos de la topología general, es equivalente al axioma de elección. Su demostración requiere lemas auxiliares (Alexander subbase theorem) y está fuera del alcance elemental de este capítulo.

\begin{exercise}
Demuestre que la topología cofinita sobre $\mathbb{N}$ es compacta pero no Hausdorff.
\end{exercise}

\begin{exercise}
Pruebe que un espacio es Hausdorff si y solo si la diagonal $\Delta=\{(x,x):\,x\in X\}$ es cerrada en $X\times X$.
\end{exercise}

\begin{exercise}
Demuestre el lema de Urysohn en espacios métricos construyendo $f(x)=d(x,A)/(d(x,A)+d(x,B))$.
\end{exercise}

\begin{problem}
Demuestre que todo espacio compacto Hausdorff es normal.
\end{problem}
